%! Logarithmic Expressions — Unit 5, Lesson 5.1 — Exit Ticket (Answer Key)
\documentclass[10pt]{article}
\usepackage{precalculus-article}
\usepackage{precalculus-key}   % pulls in -boxes; do NOT also load -boxes
\newcommand{\TallMath}[1]{$\displaystyle #1\rule[-1.4em]{0pt}{3.2em}$}

\begin{document}

\pageheader{Unit 5, Lesson 5.1}{Exit Ticket --- Answer Key}
\namedateperiod

\vspace{0.16in}

\begin{enumerate}[itemsep=20pt]

\item Convert to logarithmic form: $3^5 = 243$.

      Logarithmic form: \ans{$\log_3 243 = 5$}

      \vspace{0.04in}
      Then evaluate: $\log_3 243 = $ \ans{5}

\item Evaluate without a calculator.

      $\log_{10} 10{,}000 = $ \ans{4} \qquad\qquad $\ln e^3 = $ \ans{3}

\item Between which two consecutive integers does $\log_2 50$ fall?

      \vspace{0.04in}
      $2^{\ans{5}} = $ \ans{32} $< 50 <$ \ans{64} $= 2^{\ans{6}}$

      \vspace{0.04in}
      So $\log_2 50$ is between \ans{5} and \ans{6}.

      \vspace{0.04in}
      Explanation: \ansline{Since $2^5 = 32 < 50 < 64 = 2^6$, the exponent needed to produce 50 is between 5 and 6.}

\end{enumerate}

\end{document}
